\documentclass[11pt]{article}
\usepackage[margin=0.9in]{geometry}
\usepackage{amsmath,amssymb,booktabs,array,tabularx}
\usepackage[T1]{fontenc}
\usepackage{microtype,xcolor,tcolorbox}
\definecolor{Dynamic}{HTML}{155A83}
\definecolor{Algebraic}{HTML}{974707}
\definecolor{Link}{HTML}{284E6B}
\usepackage[colorlinks=true,linkcolor=Link,urlcolor=Link,
 citecolor=Link,filecolor=Link]{hyperref}
\newcommand{\code}[1]{\texttt{\detokenize{#1}}}
\newcommand{\dd}{\mathrm d}
\newcommand{\area}{\mathcal A}
\newcommand{\dmf}{\Delta m_{c,f}}
\newcommand{\negds}{(-\partial_r s)}
\newcommand{\buo}{\left(-\frac{\partial_r P}{\rho}\right)}
\renewcommand{\tabularxcolumn}[1]{>{\raggedright\arraybackslash}p{#1}}
\newenvironment{solved}[2]{%
 \begin{tcolorbox}[colback=#1!3!white,colframe=#1,boxrule=0.4pt,
 leftrule=2pt,arc=0pt,left=7pt,right=7pt,top=5pt,bottom=5pt,
 before skip=7pt,after skip=7pt]
 \noindent\textcolor{#1}{\textbf{#2}}\par
 \setlength{\abovedisplayskip}{5pt}\setlength{\belowdisplayskip}{5pt}
}{\end{tcolorbox}}
\newcommand{\linear}{\par\smallskip\noindent\textcolor{Dynamic}{\textbf{Static linearization.}}\par}
\setlength{\parindent}{0pt}
\setlength{\parskip}{5pt}
\setlength{\emergencystretch}{1em}
\clubpenalty=10000
\widowpenalty=10000
\numberwithin{equation}{section}
\title{RSP3 in MESA/star\\Local Entropy Moments in RSP2}
\author{}
\date{2026-09-19}
\begin{document}
\maketitle
\vspace{-1.5em}

\section{Scope and variables}
RSP3 denotes the optional three equation convection model within RSP2.
It evolves turbulent kinetic energy, entropy flux and entropy variance.
The existing independent superadiabaticity and luminosity equations remain.
The mode is selected in \code{&controls}:
\begin{verbatim}
RSP2_use_3equation_model = .true.
RSP2_alfa_pi = 1d0
RSP2_alfa_phi = 1d0
\end{verbatim}
RSP2 itself must already be enabled. The new mode defaults to false.
There is no separate RSP3 solver or \code{star_job} action.

\begin{center}
\small
\begin{tabularx}{\linewidth}{@{}l l X@{}}
\toprule
Variable & Location & Definition and role\\
\midrule
$w$ & Cell & $e_t=w^2=\tfrac12\langle|\boldsymbol v'|^2\rangle$;
 $w\geq0$, in cm\,s$^{-1}$.\\
$\Pi$ (\code{Pi}) & Face & $\langle v_r's'\rangle$, signed entropy flux,
 in erg\,cm\,g$^{-1}$\,K$^{-1}$\,s$^{-1}$.\\
$\Phi$ (\code{Phi}) & Face & $\langle s'^2\rangle$, full nonnegative entropy
 variance, in erg$^2$\,g$^{-2}$\,K$^{-2}$.\\
$Y$ & Face & $\nabla_T-\nabla_L$, signed independent solver variable.\\
$L$ & Face & Total luminosity, in erg\,s$^{-1}$.\\
$v$ or $u$ & Face or cell & Mean radial velocity; distinct from turbulent $w$.\\
\bottomrule
\end{tabularx}
\end{center}
Specific entropy is $s$. Thermodynamic derivatives are
$\chi_\rho=(\partial\ln P/\partial\ln\rho)_T$ and
$\chi_T=(\partial\ln P/\partial\ln T)_\rho$.
The full variance here is twice the $\Phi$ used by Flaskamp and Braun et al.
The old \code{PII} has entropy units and factors out turbulent velocity;
it is not the new $\Pi$.

MESA indexes inward: cell $k$ lies between outer face $k$ and inner face
$k+1$, with $k=1$ at the surface. At the face between cells $k-1$ and $k$,
\begin{equation}
 \area_f=4\pi r_f^2,\quad
 \dmf=\frac{\Delta m_{k-1}+\Delta m_k}{2},\quad
 e_{t,f}=a_fw_k^2+(1-a_f)w_{k-1}^2 .
\end{equation}
The weight is $a_f=\Delta m_{k-1}/(\Delta m_{k-1}+\Delta m_k)$ for mass
interpolation, or $1/2$ otherwise. The interpolated speed
$\langle w\rangle_f=a_fw_k+(1-a_f)w_{k-1}$ differs from $\sqrt{e_{t,f}}$.
Face EOS quantities use the selected reconstruction or RSP2 interpolation.

\textcolor{Dynamic}{Blue boxes} contain time dependent rows.
\textcolor{Algebraic}{Brown boxes} contain algebraic rows.
Unboxed formulas define their coefficients, discretization and scaling.
The physical rows are followed by their residuals and static linearizations.
Perturbations vary as $\exp(\sigma t)$; positive $\Re(\sigma)$ means growth.

\clearpage
\section{Cell turbulent energy}
The existing RSP2 energy equation remains the first turbulence equation.
The three equation mode changes its buoyant source and removes the direct
radiative sink $D_r$:
\begin{solved}{Dynamic}{Turbulent energy and finite timestep residual}
\begin{align}
 \dot e_{t,k}+P_{t,k}\frac{\dd}{\dd t}\frac1{\rho_k}
   &=-\frac{L_{t,k}-L_{t,k+1}}{\Delta m_k}+S_k-D_k+Eq_k ,
   \label{eq:et}\\
 R_{E,k}={}&w_k^2-w_{k,\mathrm{start}}^2
 +P_{t,k}^{\theta}\left(\frac1{\rho_k}-\frac1{\rho_{k,\mathrm{start}}}\right)
 \notag\\
 &+\Delta t\left[
 \frac{L_{t,k}^{\theta}-L_{t,k+1}^{\theta}}{\Delta m_k}
 -S_k+D_k-Eq_k\right]=0 .
 \label{eq:energy-residual}
\end{align}
\linear
On a static background, $\delta Eq=0$ and
\begin{equation}
 \sigma\left(2w_k\delta w_k-\frac{P_{t,k}}{\rho_k}\delta\ln\rho_k\right)
 =-\frac{\delta L_{t,k}-\delta L_{t,k+1}}{\Delta m_k}
   +\delta S_k-\delta D_k .
\end{equation}
\end{solved}
The superscript $\theta$ denotes the existing RSP2 pressure or luminosity
time weight, applied separately to each quantity. The source, dissipation
and stress heating use their existing implicit evaluation.
Adding the moments does not change these time weights.

The turbulent pressure, dissipation and kinetic energy luminosity are
\begin{align}
 P_{t,k}&=\frac23\,\code{RSP2_alfap}\,\rho_kw_k^2,\\
 D_k&=\code{RSP2_alfad}\,\frac83\sqrt{\frac23}\,
                   \frac{w_k^3}{\Lambda_k},\\
 L_{t,f}&=-\code{RSP2_alfat}\,\area_f^2
   \langle\rho^2\rangle_f\Lambda_f\langle w\rangle_f
       \frac{w_{k-1}^2-w_k^2}{\dmf}.
 \label{eq:Lt}
\end{align}
Here $\langle\rho^2\rangle_f$ is the interpolated squared density.
The physical surface and inner boundary have zero $L_t$.
Existing forced nonturbulent masks also apply.

The face pressure acceleration and cell source are
\begin{align}
 \buo_f&=\frac{\area_f(P_k-P_{k-1})}{\dmf},\\
 S_k&=\frac12\sum_{f=k,k+1}
 \buo_f\,\frac{\chi_{T,f}}{\chi_{\rho,f}C_{P,f}}\,
 \frac{w_k\Pi_f}{\sqrt{e_{t,f}}} .
 \label{eq:source}
\end{align}
The factor $w_k/\sqrt{e_{t,f}}$ reconstructs the correlation for the cell
whose energy is evolved. It makes its buoyant source vanish when $w_k=0$,
even beside a turbulent neighbor. Forced faces contribute zero.
A face with $e_{t,f}=\Pi_f=0$ contributes zero without division.
Nonzero $\Pi_f$ at zero $e_{t,f}$ is rejected as inconsistent state.
The positive energy linearization differentiates the entire reconstruction.

\textit{Source:} \code{hydro_rsp2:do1_turbulent_energy_eqn},
\code{rsp2_moment_source}, \code{compute_D}, \code{compute_Lt}.
Section~\ref{sec:numerics} specifies the residual treatment at zero $w$.

\clearpage
\section{Face entropy flux}
The covariance $\Pi=\langle v_r's'\rangle$ already includes velocity.
Its evolution contains entropy gradient driving, buoyancy acting on entropy
variance, turnover relaxation, radiative cooling and mean radial strain:
\begin{solved}{Dynamic}{Entropy flux}
\begin{equation}
\begin{aligned}
 \dot\Pi_f={}&\frac23 e_{t,f}\negds_f
 +\buo_f\frac{\chi_{T,f}}{\chi_{\rho,f}C_{P,f}}\Phi_f\\
 &-\left(\alpha_\Pi\frac{\sqrt{e_{t,f}}}{\Lambda_f}
       +\tau_{\mathrm{rad},f}^{-1}
       +(\partial_r v_{\mathrm{rad}})_f\right)\Pi_f .
\end{aligned}
\label{eq:Pi}
\end{equation}
The row actually solved is
\begin{equation}
\begin{aligned}
 R_{\Pi,f}=\frac{1}{\Pi_{\mathrm{scale},f}}\bigg\{
 &\Pi_f-\Pi_{f,\mathrm{start}}
 -\Delta t\bigg[\frac23e_{t,f}\negds_f
 +\buo_f\frac{\chi_{T,f}}{\chi_{\rho,f}C_{P,f}}\Phi_f\\
 &-\left(\alpha_\Pi\frac{\sqrt{e_{t,f}}}{\Lambda_f}
       +\tau_{\mathrm{rad},f}^{-1}
       +(\partial_r v_{\mathrm{rad}})_f\right)\Pi_f\bigg]\bigg\}=0 .
\end{aligned}
\end{equation}
\linear
\begin{equation}
\begin{aligned}
 \sigma\delta\Pi_f={}&
 \frac23\left[\negds_f\,\delta e_{t,f}
                      +e_{t,f}\,\delta\negds_f\right]\\
 &+\buo_f\frac{\chi_{T,f}}{\chi_{\rho,f}C_{P,f}}\,\delta\Phi_f
 +\Phi_f\,\delta\left[
 \buo\frac{\chi_T}{\chi_\rho C_P}\right]_f\\
 &-\left(\alpha_\Pi\frac{\sqrt{e_{t,f}}}{\Lambda_f}
                 +\tau_{\mathrm{rad},f}^{-1}\right)\delta\Pi_f\\
 &-\Pi_f\,\delta\left[
 \alpha_\Pi\frac{\sqrt{e_{t,f}}}{\Lambda_f}
 +\tau_{\mathrm{rad},f}^{-1}
 +(\partial_r v_{\mathrm{rad}})_f\right].
\end{aligned}
\end{equation}
\end{solved}
The static mean strain is zero, but its perturbation remains.
For the two velocity grids, the implemented face strain is
\begin{align}
 (\partial_r u)_f&=
       \frac{\area_f\rho_f}{\dmf}(u_{k-1}-u_k),
       &&\code{u_flag},\\
 (\partial_r v)_f&=
 a_f\frac{v_k-v_{k+1}}{r_k-r_{k+1}}
 +(1-a_f)\frac{v_{k-1}-v_k}{r_{k-1}-r_k},
       &&\code{v_flag}.
\end{align}
The face velocity expression reproduces an affine radial velocity on an
unequal grid. The strain term is a deformation of the covariance;
it introduces no new eddy viscosity coefficient.

At fixed mass weights,
$\delta e_{t,f}=2a_fw_k\delta w_k+
2(1-a_f)w_{k-1}\delta w_{k-1}$.
The Jacobian retains derivatives of face EOS, pressure differences,
mixing length and entropy driving. The accepted state and row scale are
held fixed in a Newton solve.

\textit{Source:} \code{hydro_rsp2:rsp2_moment_rhs},
\code{rsp2_buoyancy_face}, \code{do1_rsp2_moment_eqns}.

\clearpage
\section{Face entropy variance and coefficients}
\begin{solved}{Dynamic}{Full entropy variance}
\begin{align}
 \dot\Phi_f&=2\negds_f\Pi_f-
 \left(\alpha_\Phi\frac{\sqrt{e_{t,f}}}{\Lambda_f}
                    +2\tau_{\mathrm{rad},f}^{-1}\right)\Phi_f ,
 \label{eq:Phi}\\
 R_{\Phi,f}&=\frac{\Phi_f-\Phi_{f,\mathrm{start}}
 -\Delta t\left[
 2\negds_f\Pi_f-
 \left(\alpha_\Phi\frac{\sqrt{e_{t,f}}}{\Lambda_f}
       +2\tau_{\mathrm{rad},f}^{-1}\right)\Phi_f\right]}
 {\Phi_{\mathrm{scale},f}}=0 .
\end{align}
\linear
\begin{equation}
\begin{aligned}
 \sigma\delta\Phi_f={}&2\negds_f\,\delta\Pi_f+
                      2\Pi_f\,\delta\negds_f\\
 &-\left(\alpha_\Phi\frac{\sqrt{e_{t,f}}}{\Lambda_f}
       +2\tau_{\mathrm{rad},f}^{-1}\right)\delta\Phi_f\\
 &-\Phi_f\,\delta\left[
 \alpha_\Phi\frac{\sqrt{e_{t,f}}}{\Lambda_f}
       +2\tau_{\mathrm{rad},f}^{-1}\right].
\end{aligned}
\end{equation}
\end{solved}
There is no spatial transport term in either moment equation. Nonzero
\code{RSP2_alfat} transports kinetic energy through $L_t$ only.
The full variance convention doubles its driving relative to a half
variance; it does not change the turnover decay rate.

The radiative rate is
\begin{equation}
 \tau_{\mathrm{rad},f}^{-1}=
 \frac{4\sigma_{\mathrm{SB}}\gamma_r^2T_f^3}
 {C_{P,f}\kappa_f\rho_f^2\Lambda_f^2},
 \qquad \gamma_r=2\sqrt3\,\code{RSP2_alfar}.
 \label{eq:cooling}
\end{equation}
Thus \code{RSP2_alfar = 0} is permitted. It removes radiative damping
of both moments while retaining turnover relaxation.
There is no direct $D_r$ in Equation~\eqref{eq:et}.
The one equation model instead applies $D_r=e_t/\tau_{\rm rad}$ in cells.

\begin{center}
\small
\begin{tabularx}{\linewidth}{@{}l l X@{}}
\toprule
Control & Default & Coefficient in this formulation\\
\midrule
\code{RSP2_alfa_pi} & 1 & $\alpha_\Pi=6\sqrt{2/3}\,\code{RSP2_alfa_pi}$\\
\code{RSP2_alfa_phi} & 1 & $\alpha_\Phi=4\sqrt{2/3}\,\code{RSP2_alfa_phi}$\\
\code{RSP2_alfad} & 1 & $C_D=(8/3)\sqrt{2/3}\,\code{RSP2_alfad}$\\
\code{RSP2_alfar} & 0 & $\gamma_r=2\sqrt3\,\code{RSP2_alfar}$\\
\code{RSP2_alfat} & 0 & $\alpha_\omega=\code{RSP2_alfat}$\\
\code{RSP2_alfap} & 0 & Multiplier of $(2/3)\rho e_t$\\
\code{RSP2_alfam} & 0.25 & Existing eddy viscosity coefficient\\
\bottomrule
\end{tabularx}
\end{center}
The unit Pi/Phi controls give the local MLT coefficients listed by
Braun et al.~\cite{braun}, $\alpha_\Pi\simeq4.899$ and
$\alpha_\Phi\simeq3.266$. In their fiducial solar model,
$\alpha_\Pi=2.155$ and $\alpha_\Phi=2$ multiply spatial transport terms.
They are not calibrations of the local decay terms used here.
Their case A replaces both closures only in the outer layers and uses
$\alpha_\Pi=2.88$, $\alpha_\Phi=2$ there; the interior remains nonlocal.
Neither case supplies a calibration for this fully local pulsation model.
The radiative and kinetic transport normalizations do retain their meanings:
$\gamma_r\simeq3.46$ corresponds to \code{RSP2_alfar = 1}, and
$\alpha_\omega=0.25$ to \code{RSP2_alfat = 0.25}.
The two moment controls must be finite and positive.
The mode requires \code{RSP2_source_seed = 0}.

\clearpage
\section{Luminosity and temperature equations}
The new entropy flux sets convective luminosity directly:
\begin{align}
 L_{c,f}&=\area_f\rho_fT_f\Pi_f,
 &\delta L_{c,f}
 &=\area_f\rho_fT_f\,\delta\Pi_f+
                  \Pi_f\,\delta(\area_f\rho_fT_f),\\
 L_{r,f}&=\mathrm{Lrad\_coeff}_f(\nabla_{L,f}+Y_f),
 &\mathrm{Lrad\_coeff}_f&=
 \area_f\frac{4acT_f^4}{3\kappa_f\rho_fH_{P,f}} .
\end{align}
Here $a$ is the radiation constant. There is no extra $w$ multiplying
$\Pi$ and no enthalpy flux limiter. $L_t$ is Equation~\eqref{eq:Lt}.
\begin{solved}{Algebraic}{Independent flux balance at interior faces}
\begin{align}
 R_{\mathrm{flux},f}&=
 \frac{L_{r,f}+L_{c,f}+L_{t,f}-L_f}{L_{\mathrm{scale},f}}=0,\\
 0&=\delta L_{r,f}+\delta L_{c,f}+\delta L_{t,f}-\delta L_f,\\
 \delta L_{r,f}&=\mathrm{Lrad\_coeff}_f(\delta\nabla_{L,f}+\delta Y_f)
             +(\nabla_{L,f}+Y_f)\,\delta\mathrm{Lrad\_coeff}_f .
\end{align}
\end{solved}
The temperature row is retained in \code{i_equL}. The separate flux
row is \code{i_rsp2_flux}; it supplies the extra equation for $Y$.
Using $Y$ as a solver variable avoids locally inverting a saturated flux law.
It does not by itself prove stability or realizability of the moment closure.

\begin{solved}{Algebraic}{Supported interior temperature residuals}
The default QHSE row is
\begin{equation}
 R_T=\dmf\left(\frac{\dd\ln P}{\dd m}\right)_{\rm QHSE}
       (\nabla_L+Y)-\frac{T_{k-1}-T_k}{T_*}=0.
 \label{eq:qhs}
\end{equation}
Here $T_*=(\Delta m_{k-1}T_k+\Delta m_kT_{k-1})/
(\Delta m_{k-1}+\Delta m_k)$.
The radiation pressure row is, up to its positive start normalization,
\begin{equation}
 R_T=-\frac{\Delta m_{T,f}\kappa_{\rm row,f}L_{r,f}}
                   {c\area_f^2\lambda_f}
               -(P_{{\rm rad},k-1}-P_{{\rm rad},k})=0,
 \quad P_{\rm rad}=\frac{aT^4}{3}.
 \label{eq:dprad}
\end{equation}
The actual logarithmic gradient row is
\begin{equation}
 R_T=(\nabla_L+Y)(\ln P_{k-1}-\ln P_k)
                       -(\ln T_{k-1}-\ln T_k)=0 .
 \label{eq:actual}
\end{equation}
Their linearized rows are $\delta R_T=0$, with current coefficients,
face reconstruction and $\delta\nabla_T=\delta\nabla_L+\delta Y$
differentiated throughout.
\end{solved}
In Equation~\eqref{eq:dprad}, $\Delta m_{T,f}$ includes the native boundary
spacing choice, $\kappa_{\rm row}$ the opacity floor, and $\lambda$ the
existing radiative flux factor. This radiation treatment is independent of
the removed enthalpy limiter. The QHSE coefficient retains the selected
momentum geometry and time weights.

At $k=1$, the independent flux row sets $Y_1=0$.
The separate surface temperature or luminosity condition remains;
the RSP2 surface option uses $L_1=f_{\rm surf}\area_1caT_1^4$ with its
native area time weight. The nonlinear \code{constant_L} option replaces
the temperature row with $L_k=L_1$ and is excluded from LNA.

\clearpage
\section{Entropy driving and dynamic gradL}
The moment equations need the thermal entropy differential represented by
the selected temperature row. For uniform composition,
\begin{equation}
 \negds_f=C_{P,f}\left[
       (-\partial_r\ln T)_{{\rm row},f}
       -\nabla_{{\rm ad},f}(-\partial_r\ln P)_{{\rm actual},f}\right].
 \label{eq:entropy}
\end{equation}
Using a logarithmic entropy difference with a nonlogarithmic temperature
row would shift the discrete neutral state. The implementation therefore
maps the actual row before forming the small superadiabatic difference.

For the QHSE and radiation pressure rows, the two thermal coefficients are
\begin{align}
 (-\partial_r\ln T)_{\nabla_T=1,f}
 &=\begin{cases}
 -\area_f\rho_f(\dd\ln P/\dd m)_{\rm QHSE}\,T_*/T_f,
                      &\text{QHSE},\\[2pt]
 \displaystyle\frac{\rho_f\kappa_{\rm row,f}\mathrm{Lrad\_coeff}_f}
       {4c\area_f\lambda_fP_{{\rm rad},f}},
                      &\text{radiation pressure},
 \end{cases}\\
 (-\partial_r\ln P)_{{\rm actual},f}
 &=\frac{\area_f\rho_f}{\Delta m_{T,f}}\,
                     \frac{P_k-P_{k-1}}{P_f}.
\end{align}
For QHSE, $\Delta m_{T,f}=\dmf$.
The reference pressure differential uses $P_k-P_{k-1}$ when
\path{TDC_use_dynamical_gradL} is true and hydro velocity is active.
Otherwise it uses the pressure jump implied by the QHSE momentum
coefficient, including its mass and area conventions.
Writing that choice explicitly,
\begin{align}
 (-\partial_r\ln P)_{{\rm ref},f}
 &=\frac{\area_f\rho_f}{\Delta m_{T,f}P_f}
 \begin{cases}
 P_k-P_{k-1},&\text{dynamic gradL},\\
 \Delta P_{{\rm HSE},f},&\text{static gradL},
 \end{cases}\\
 \nabla_{L,f}
 &=(\nabla_{{\rm ad},f}+B_f)
 \frac{(-\partial_r\ln P)_{{\rm ref},f}}
      {(-\partial_r\ln T)_{\nabla_T=1,f}} .
\end{align}
Here $B_f$ is the existing \code{gradL_composition_term}, included only for
the Ledoux choice. It is held fixed in thermal AD.
The evaluated driving is
\begin{equation}
\begin{aligned}
 \negds_f=C_{P,f}\big[&
 (-\partial_r\ln T)_{\nabla_T=1,f}\,Y_f
 +B_f(-\partial_r\ln P)_{{\rm ref},f}\\
 &+\nabla_{{\rm ad},f}
 \{(-\partial_r\ln P)_{{\rm ref},f}
             -(-\partial_r\ln P)_{{\rm actual},f}\}\big].
\end{aligned}
\label{eq:driving}
\end{equation}
The last difference is identically zero in the dynamic branch and is
removed algebraically before evaluation. This retains the small $Y$
contribution without subtracting two nearly equal neutral contributions.
All structure derivatives remain implicit; this changes no luminosity
time weight.

For the logarithmic row, $\nabla_L=\nabla_{\rm ad}+B$ and
\begin{equation}
 \negds_f=C_{P,f}\frac{\area_f\rho_f}{\dmf}
                   \ln(P_k/P_{k-1})(Y_f+B_f).
\end{equation}
The logarithm is evaluated using \code{log1p}.
For nonlinear \code{constant_L}, the code instead uses the measured
$C_P\area\rho/\dmf\,[(T_k-T_{k-1})/T_f-\nabla_{\rm ad}(P_k-P_{k-1})/P_f]$.
The composition term supplies the existing thermal correction; the model
does not evolve composition covariances.

\textit{Source:} \code{hydro_gradient_support:get_rsp2_thermal_gradient}.
LNA uses the same mapping with instantaneous coefficients.

\clearpage
\section{Gas energy and pressure work}
Both RSP2 energy forms already include turbulent storage.
The moment model changes $L_c$ and $S$ but adds no second gas source
$-S+D$. Combining the total and turbulent energy rows recovers that exchange.
Let $D_m L=(L_k-L_{k+1})/\Delta m_k$ and let $\epsilon$ denote the existing
nuclear, neutrino and other heating terms.
\begin{solved}{Dynamic}{The \code{dedt} energy row}
\begin{align}
 \dot e_{\rm eff}&=-D_mL+\epsilon+\epsilon_{\rm visc}
                            -\mathcal W,\\
 R_{\rm dedt}&=\epsilon+\epsilon_{\rm visc}-D_mL^\theta
 -\mathcal W^\theta-\frac{e_{\rm eff}-e_{\rm eff,start}}{\Delta t}=0,\\
 \sigma\delta e_{\rm eff}&=-D_m\delta L+\delta\epsilon-\delta\mathcal W
                    \quad\text{on a static background}.
\end{align}
\end{solved}
The finite step equation is written before the native row normalization.
The work control selects both the pressure work and stored energy:
\begin{align}
 e_{\rm eff}&=e+e_t,
 &\mathcal W_{\rm local}
 &=P_{c,k}\frac{\area_kv_{f,k}-\area_{k+1}v_{f,k+1}}{\Delta m_k},
 \label{eq:localwork}\\
 e_{\rm eff}&=e+e_t+K+E_{\rm grav},
 &\mathcal W_{\rm flux}
 &=\frac{P_{f,k}\area_kv_{f,k}
               -P_{f,k+1}\area_{k+1}v_{f,k+1}}{\Delta m_k}.
 \label{eq:fluxwork}
\end{align}
Equation~\eqref{eq:localwork} applies when
\code{use_P_d_1_div_rho_form_of_work} is true.
Equation~\eqref{eq:fluxwork} applies when it is false.
$P_c$ and $P_f$ are the native work pressures including selected turbulent
pressure. With \code{u_flag}, $v_f$ is the differentiated Riemann contact
velocity and $P_f$ is the corresponding work pressure; with \code{v_flag},
$v_f$ is the stored face velocity. The native discrete $K$ and
$E_{\rm grav}$ are retained.

\begin{center}
\small
\begin{tabular}{@{}lll@{}}
\toprule
Grid & Local work: $\epsilon_{\rm visc}$ & Flux work: $\epsilon_{\rm visc}$\\
\midrule
$v$ & $Eq_k$ &
$Eq_k+\tfrac12(\bar v_kUq_k+\bar v_{k+1}Uq_{k+1})$\\
$u$ & $Eq_k$ & $Eq_k+\bar u_kUq_k$\\
\bottomrule
\end{tabular}
\end{center}
Bars denote accepted/current averages. The nonlinear $v$ power retains
the cell mass correction when selected. These viscous powers are second
order about a static model.

\begin{solved}{Dynamic}{The RSP2 \code{eps_grav} energy row}
\begin{align}
 R_{\epsilon_{\rm grav}}={}&\epsilon+\epsilon_{\rm grav}+Eq-D_mL^\theta
 -\mathcal W_{\neg{\rm EOS}}^\theta
 -\frac{e_t-e_{t,\mathrm{start}}}{\Delta t}=0,\\
 \mathcal W_{\neg{\rm EOS}}
 &=(P_c-P_{\rm EOS})
 \frac{\area_kv_{f,k}-\area_{k+1}v_{f,k+1}}{\Delta m_k}.
\end{align}
\end{solved}
This option always uses local work, with no explicit $K$, $E_{\rm grav}$
or $vUq$ storage/power branch. At fixed composition and consistent EOS,
$\epsilon_{\rm grav}=-\dot e-P_{\rm EOS}\,\dd(1/\rho)/\dd t$.
Its static inertia, for the ordinary thermal coefficients, is
$2w\delta w+TC_P[(1-\nabla_{\rm ad}\chi_T)\delta\ln T
-\nabla_{\rm ad}\chi_\rho\delta\ln\rho]$.
Native latent heat, entropy blending and energy factors retain their
existing treatment. Adding EOS pressure work again would double count it.

\clearpage
\section{Eddy viscosity on the two velocity grids}
The new moment rows retain the RSP2 $Eq$ and $Uq$ discretization.
TDC uses separate hydro viscosity routines because its turbulent velocity
is a face quantity. Only the public length functions are shared.
The following formulas use native masses without optional mass corrections.

\textbf{Face mean velocity $v$: cell stress.}
\begin{align}
 Chi_k&=\frac{16\pi}{3}\frac{\alpha_m\rho_k^2\Lambda_kw_k}{\Delta m_k}
       \frac{r_k^6+r_{k+1}^6}{2}
       \left(\frac{v_k}{r_k}-\frac{v_{k+1}}{r_{k+1}}\right),\\
 Uq_k&=\frac{4\pi}{r_k^{\rm work}\Delta m_{\rm face,k}}
                          (Chi_{k-1}-Chi_k),\\
 Eq_k&=\frac{4\pi Chi_k}{\Delta m_k}
       \left(\frac{v_k^{\rm work}}{r_k^{\rm work}}
              -\frac{v_{k+1}^{\rm work}}{r_{k+1}^{\rm work}}\right).
\end{align}
Here $\alpha_m=\code{RSP2_alfam}$,
$\Delta m_{\rm face,k}=(\Delta m_{k-1}+\Delta m_k)/2$ internally,
and $\Delta m_{\rm face,1}=\Delta m_1/2$.
The stress uses current geometry and velocity; the force and work use the
native optional time centered radii. The exterior stress is zero.

\textbf{Cell mean velocity $u$: face stress.}
\begin{align}
 Chi_{f,k}&=\frac{16\pi}{3}
 \frac{\alpha_m\rho_{f,k}^2r_k^6\Lambda_{f,k}\langle w\rangle_{f,k}}
      {\Delta m_{\rm face,k}}
 \left(\frac{u_{k-1}}{r_{\rm mid,start,k-1}}
                  -\frac{u_k}{r_{\rm mid,start,k}}\right),\\
 \Delta m_kUq_k&=\frac{4\pi}{r_{\rm mid,start,k}}
                         (Chi_{f,k}-Chi_{f,k+1}),\\
 Eq_{f,k}&=\frac{4\pi Chi_{f,k}}{\Delta m_{\rm face,k}}
 \left(\frac{u_{k-1}^{\rm work}}{r_{\rm mid,start,k-1}}
                  -\frac{u_k^{\rm work}}{r_{\rm mid,start,k}}\right),\\
 Eq_k&=\tfrac12(Eq_{f,k}+Eq_{f,k+1}).
\end{align}
The density is the native viscosity face density, while $\langle w\rangle_f$
uses the RSP2 weights. \code{compute_Uq_dm_cell} returns force:
the momentum row divides by $\Delta m_k$ once.
The half sum of face \emph{specific} heating is essential. On an unequal
mesh, weighting each cell by $\Delta m_k$ allocates each face's full dual
mass to that face heating.

For flux work, the work velocities are always
$\bar v=(v+v_{\rm start})/2$ or $\bar u=(u+u_{\rm start})/2$.
For local work they are averaged only when velocity time centering is on;
otherwise they are current velocities.
With the same stress in force and heating, summation by parts gives
\begin{align}
 \sum_k\Delta m_{\rm face,k}v_k^{\rm work}Uq_k
       +\sum_k\Delta m_kEq_k&=\text{boundary stress power},
                    &&v\text{ grid},\\
 \sum_k\Delta m_ku_k^{\rm work}Uq_k
       +\sum_k\Delta m_kEq_k&=\text{boundary stress power},
                    &&u\text{ grid}.
\end{align}
The boundary power is zero for zero boundary stress or velocity.
An excised inner boundary can retain the native optional stress and its
half cell mass. Forced regions mask stress while retaining its divergence
at the interface. A full center uses the zero stress limit.

For a static background, $\delta Chi$ contains the perturbed velocity
strain times the background coefficient; $\delta Uq$ stays in momentum.
$Eq$ and $vUq$ are second order, so they add no first order heating row.
These power identities concern residual values. The existing
block tridiagonal conservative energy Jacobian can omit a $k+2$ partial;
the moment equations do not repair that separate stencil limitation.

\clearpage
\section{Scales and the zero turbulence boundary}
\label{sec:numerics}
The independent luminosity equation uses a fixed accepted state scale,
\begin{equation}
 L_{\rm scale,f}=
 \max\left(1\ {\rm erg\,s^{-1}},\ |L_{f,\rm start}|,
                   10^{-3}\max_j|L_{j,\rm start}|\right).
\end{equation}
At initialization, the current state supplies this scale until an accepted
start state exists. Moment row and column scales use the same luminosity
reference and accepted face geometry and EOS:
\begin{align}
 \Pi_{\rm ref,f}&=\frac{L_{\rm scale,f}}{\area_f\rho_fT_f},
 &v_{\rm ref,f}&=\left(\frac{L_{\rm scale,f}}{\area_f\rho_f}\right)^{1/3},\\
 \Pi_{\rm scale,f}&=\max(|\Pi_{f,\rm start}|,\Pi_{\rm ref,f}),
 &\Phi_{\rm scale,f}&=\max\left(\Phi_{f,\rm start},
                     \frac32\frac{\Pi_{\rm ref,f}^2}{v_{\rm ref,f}^2}\right).
\end{align}
These are numerical normalizations, not bounds on heat transport.
They are not differentiated as changing physics during a Newton solve.
The new correction weights are reciprocal to the larger of the scale and
the current variable magnitude. The signed $Y$ correction weight remains
$1/(1+|Y|)$; it neither changes sign nor caps $Y$.

In the three-equation model, outside the local zero start branch below,
the normalized energy row is
\begin{equation}
 R_w=\frac{\mathrm{energy\_eqn\_scal}}{\Delta t}
 \begin{cases}
 c_s\,[R_E/w]_{\rm analytic},&0<w<c_s,\\
 R_E,&w\geq c_s.
 \end{cases}
 \label{eq:wscale}
\end{equation}
This equals $\max(1,c_s/w)R_E\,\mathrm{energy\_eqn\_scal}/\Delta t$
at positive $w$. The one-equation RSP2 model retains the original normalized
energy residual for nonzero accepted energy. This rescaling preserves the
physical energy balance and does not
weaken its residual tolerance. Crucially, the code cancels known factors
\emph{before} differentiation; it does not divide an assembled $R_E$.

\paragraph{What cancels.}
The storage and local sinks give
\begin{align}
 \frac{w^2-w_{\rm start}^2}{w}
   &=w-\frac{w_{\rm start}^2}{w},
 &\partial_w\left(w-\frac{w_{\rm start}^2}{w}\right)
   &=1+\left(\frac{w_{\rm start}}w\right)^2,\\
 \frac{D}{w}&=\frac{C_Dw^2}{\Lambda},
 &\frac{D_r}{w}&=\frac{w}{\tau_{\rm rad}}.
\end{align}
Here $C_D=(8/3)\sqrt{2/3}\,\code{RSP2_alfad}$; $D_r=0$ in the
three equation mode. Its buoyancy source becomes
\begin{equation}
 \left[\frac{S_k}{w_k}\right]_{\rm analytic}
 =\frac12\sum_{f=k,k+1}
 \buo_f\frac{\chi_{T,f}}{\chi_{\rho,f}C_{P,f}}
 \frac{\Pi_f}{\sqrt{e_{t,f}}}.
\end{equation}
Only the explicit numerator $w_k$ in Equation~\eqref{eq:source} cancels.
The face energy still depends on both neighboring cell velocities, and its
derivatives remain. One-equation RSP2 uses the original energy residual
when the source seed is nonzero.

The uncanceled source quotient subtracts derivative terms of order
$\code{Source_div_w}/w$; their roundoff can overwhelm the storage derivative.
Cancel that factor first. For each remaining numerator $x$, evaluate
$q=x/w$ and then $\delta q=(\delta x-q\,\delta w)/w$.
This avoids forming $1/w^2$, which can overflow even when $q$ and its
derivatives are finite. No physical coefficient is frozen.

\clearpage
\paragraph{The complete divided balance.}
For positive $w_k$, the implementation evaluates
\begin{align}
 [R_{E,k}/w_k]_{\rm analytic}={}&w_k
 +\theta_P\frac23\code{RSP2_alfap}\,\rho_kw_k\Delta(1/\rho_k)
 \notag\\
 &+\frac{(1-\theta_P)P_{t,k,\rm start}\Delta(1/\rho_k)
                -w_{k,\rm start}^2}{w_k}\notag\\
 &+\Delta t\biggl[
 \frac{[L_{t,k}^{\theta}/w_k]_{\rm analytic}
       -[L_{t,k+1}^{\theta}/w_k]_{\rm analytic}}{\Delta m_k}
 \notag\\
 &\hspace{3em}-[S_k/w_k]_{\rm analytic}+\frac{C_Dw_k^2}{\Lambda_k}
 +\frac{w_k}{\tau_{\rm rad,k}}-[Eq_k/w_k]_{\rm analytic}\biggr].
 \label{eq:factored-energy}
\end{align}
Here $\Delta(1/\rho_k)=1/\rho_k-1/\rho_{k,\rm start}$ and the radiative
term is omitted in the three equation mode. Accepted energy and pressure
remain in the second line; they cannot be canceled with current $w_k$.

\paragraph{Face transport and viscosity.}
$L_t$ and $u$ grid heating depend on neighboring turbulence and need not
vanish at $w_k=0$. For face $f$ between cells $k-1,k$,
\begin{equation}
 \frac{\langle w\rangle_f}{w_k}
   =a_f+(1-a_f)\frac{w_{k-1}}{w_k},\qquad
 \frac{\langle w\rangle_f}{w_{k-1}}
   =a_f\frac{w_k}{w_{k-1}}+(1-a_f).
\end{equation}
Only the interpolated velocity factor in $L_t$ and face stress is replaced;
energy jumps, mass weights, strain and all derivatives remain.
The $u$ grid retains half of each bounding face's specific heating.
At an active inner boundary,
$\langle w\rangle/w_N=w_N/w_N=1$ is canceled exactly.
On the $v$ grid, cell stress and heating already have a common local $w$;
the existing \code{Eq_div_w} supplies their divided form.

The luminosity time weight is retained explicitly:
\begin{equation}
 [L_{t,f}^{\theta}/w_k]_{\rm analytic}
 =\theta_L[L_{t,f}/w_k]_{\rm analytic}
 +(1-\theta_L)\frac{L_{t,f,\rm start}}{w_k}.
\end{equation}
This retains nonzero $\alpha_t$, both velocity grids and both work forms.
Physical $L_t$, $Eq$, $Uq$ and the continuous LNA equations are unchanged.

\paragraph{The zero state.}

When $e_{t,\rm start}=0$, $\alpha_t=0$, the source seed is zero, and either
$v$ is used or $\alpha_m=0$, the local row has a common factor $w$.
The implementation cancels it analytically and solves the branch equation
\begin{equation}
 \frac{c_s\,\mathrm{energy\_eqn\_scal}}{\Delta t}
       \min\left(w,\left[R_E/w\right]_{\rm analytic}\right)=0.
\end{equation}
The accepted energy test includes exact underflow of $w_{\rm start}^2$
to zero; no positive representable energy is discarded.
The dormant branch wins an equality tie; no small denominator is inserted.
For a nonlocal row with $e_{t,\rm start}=w=R_E=0$, the dormant row is
$w c_s\,\mathrm{energy\_eqn\_scal}/\Delta t=0$.
If transport or heating supplies finite energy, the predictor supplies
a positive trial state before solving the original energy balance.

Accepted $w$ and $\Phi$ are nonnegative; $\Pi$ remains signed.
An active positive $w$ correction leaves at least one tenth of its current
value; local branch and forced cells may reach zero. No cached negative
$w$ substitution is used.

The predictor uses two directional sweeps and a positive quadratic storage
estimate, including source, damping, $Eq$, accepted/current $L_t$, pressure
work and surviving $\Phi$ at negligible face energy. The existing buoyancy
guess also applies when both adjacent $w$ are below machine precision times
$\Delta t\,|\buo_f\chi_{T,f}/(\chi_{\rho,f}C_{P,f})|\sqrt{\Phi_f}$.
It changes trial guesses only. Accepted history and time weights remain
fixed; timestep convergence still requires testing.

\clearpage
\section{Initialization, lengths and mesh state}
\textbf{Enabling the mode.}
The internal flag records the loaded layout. Enabling the mode adds
$\Pi,\Phi$ after $Y$, preserves the other hydro variables and initializes
missing moments. Compute the target neutral reference before assigning
and unpacking
\begin{equation}
 Y_{f,0}=\nabla_{T,f,\rm old}-\nabla_{L,f,\rm new}.
\end{equation}
For an unforced face with positive entropy driving
$-\partial s/\partial r>0$, positive old convective luminosity and
positive face energy, the missing moments are initialized as
\begin{equation}
 \Pi_{f,0}=\frac{L_{c,f,\rm old}}{\area_f\rho_fT_f},
 \qquad
 \Phi_{f,0}=\frac{3\Pi_{f,0}^2}{2e_{t,f}}.
\end{equation}
These faces retain their heat flux with a fully correlated parcel variance.
Other faces start with $\Pi=\Phi=0$, retaining cell turbulent energy.
This changes their imported flux and avoids inferring finite variance
from tiny inherited velocities in stable layers. The entropy gradient
generates moments when turbulent energy is present. This initial state
does not imply moment equilibrium. Retained nonzero flux at zero face
energy is rejected. Surface $Y$ is reset to zero.

Conversion sets current and accepted moments consistently and invalidates
unavailable earlier generations. It precedes step snapshots and bypasses
existing three equation restarts, Newton iterations and retries.
Accepted moments use the native \code{xh_start} state.

\textbf{Boundary and dormant rows.}
Moments vanish at faces $k=1,N$, for zero mixing length parameter, and
in the selected outer and inner forced nonturbulent regions. Their rows are
\begin{equation}
 R_{\Pi,f}=\Pi_f/\Pi_{\rm scale,f}=0,\qquad
 R_{\Phi,f}=\Phi_f/\Phi_{\rm scale,f}=0 .
\end{equation}
A face with both adjacent energies and current and accepted moments zero
also uses this branch. Finite adjacent turbulent energy retains the
moment evolution equations.

\textbf{Public TDC length functions.}
RSP2 calls \code{get_TDC_Hp_face},
\code{get_TDC_mixing_length_face} and \code{get_TDC_mixing_length_cell}.
The lengths are analytic. Cell lengths use the cell EOS and gravity
averaged over bounding faces. For the HSE choice,
\begin{equation}
 H_{P,k}=\frac{P_k}{\rho_k(g_k+g_{k+1})/2},\qquad
 \Lambda_{0,k}=\alpha_{\rm MLT}H_{P,k}.
\end{equation}
The shared harmonic option evaluates
$\Lambda_k=(\Lambda_{0,k}^{-1}+(\beta_h r_k^{\rm cell})^{-1})^{-1}$
when $\beta_h>0$, with the native cell radius. The same helpers support
the other length choices and face reconstruction.

\begin{center}
\small
\begin{tabularx}{\linewidth}{@{}l X X@{}}
\toprule
Mesh path & Cell $w^2$ & Face $Y,\Pi/w_f,\Phi$\\
\midrule
Normal remesh & Conservative energy remap &
Existing monotonic cubic face interpolation, \code{interp_pm}.\\
Envelope remesh & Conservative energy remap &
Existing envelope face interpolation.\\
Split and merge & Preserve cell integrated turbulent energy &
Split: interpolate a new interior face in mass; merge: retain surviving values of these quantities.\\
\bottomrule
\end{tabularx}
\end{center}
Moments are face point values. All three paths first retain the old ratio
$\Pi_f/w_f$, with $w_f=\sqrt{a_fw_k^2+(1-a_f)w_{k-1}^2}$, and take
the ratio as zero at zero old face energy. After the conservative cell
energy remap, they reconstruct
\begin{equation}
 \Pi_{f,\mathrm{new}}=w_{f,\mathrm{new}}\,
 \mathcal I_f\!\left[\frac{\Pi_{\mathrm{old}}}{w_{\mathrm{old}}}\right].
\end{equation}
Here $\mathcal I_f$ is the face interpolation in the table. This prevents
a new face inside a quiet cell from inheriting a finite boundary flux
without the turbulent energy that carried it. Surviving faces in split
and merge also use the new face energy. If $w_{f,\mathrm{new}}=0$, Pi is
zero while Phi retains its entropy variance reservoir. This mesh transfer
does not cap Pi at positive energy or change its evolution equation.
Remeshing reapplies forced masks and rejects negative variance.

The ordinary RSP2 thermal remap reconstructs cell averages with the native
monotonic polynomials. With \code{mesh_adjust_get_T_from_E} false it
integrates $\ln T$; with the control true it integrates specific internal
energy and uses the existing EOS inversion and optional gravitational
and kinetic energy correction. That correction uses the same polynomial
reconstruction of old specific gravitational plus kinetic energy; its
existing relative cap is retained. The polynomial constant retains the old
cell average after curvature limiting. At an excised inner boundary the
last cell uses a one-sided slope; the actual center retains zero slope.
The energy slope is bounded to keep its endpoint values nonnegative.
Temperature remapping does not itself conserve internal energy.

Photo version 21 stores the mode and moments; version 20
without moments remains readable. Saved models use bit 17 and Pi/Phi columns.
Retries and old models use the generic hydro state machinery.

\section{Linear nonadiabatic operator and physical limits}
The radial LNA system is
\begin{equation}
 \mathbf A\,\delta\mathbf x=\sigma\mathbf B\,\delta\mathbf x .
\end{equation}
The RSP2 variables $w,Y$ are supplemented by face $\Pi,\Phi$.
Each active moment has unit diagonal inertia in $\mathbf B$.
The corresponding $\mathbf A$ row is the AD derivative of the continuous
right hand side in Equations~\eqref{eq:Pi} and \eqref{eq:Phi}.
No finite timestep, Newton scale or predictor enters this operator.
The entropy mapping uses instantaneous coefficients, including dynamic
gradL when selected.

The turbulent energy row is the continuous linearization on page 2.
The luminosity row includes $\delta L_c=\delta(\area\rho T\Pi)$.
The chosen frozen convective flux option instead holds that luminosity
perturbation fixed, while retaining the moment equations.
The two cases therefore define different operators.

Forced or fully dormant moment rows impose
$\delta\Pi=\delta\Phi=0$ with zero inertia.
A fully dormant cell similarly uses $\delta w=0$. At zero represented
cell energy, including underflow of $w^2$, a local row without a source
seed or kinetic transport also uses this constraint when
$S_k/w_k\leq0$. The source ratio is evaluated in its factored form.
This is the zero branch of $\min(w_k,R_{E,k}/w_k)=0$ at a static background;
it uses current energy, not accepted timestep history.
Positive driving or nonlocal coupling retains the physical energy row.
An interface with finite face energy retains the dynamic moment rows,
even when an adjacent cell has $\delta w=0$.
Active moments at exactly zero face energy have no regular linearization
in these variables; LNA reports an error for that background.
This is a linearization of the selected dormant branch and does not
describe spontaneous onset from an exactly empty turbulent state.
A temperature selected LNA cut retains the original physical masks;
it does not create an extra forced moment boundary.

On a static background, pressure work retains the background pressure times
the perturbed face volume flux. For \code{dedt} with flux work,
$\delta K=0$ but $\delta E_{\rm grav}$ remains.
For local work and \code{eps_grav}, mechanical inertia is absent.
The first order viscous force remains while $\delta Eq=0$.
These distinctions are detailed for all supported hydro branches in the
companion \href{star_LNA_manuscript.pdf}{star LNA manuscript}.

\textbf{Equilibrium and diagnostics.}
LNA assumes a stationary background of the selected equations.
Enabling the mode with the prescribed initial moments does not establish that
condition. Setup reports the dimensional maxima of
$|\dot\Pi_f|$ and $|\dot\Phi_f|$, so a residual thermal moment drift is
visible before interpreting growth rates.
The eigenfunction output includes real and imaginary Pi/Phi components.
Profiles expose \code{Pi}, \code{Phi} and \code{Pi_covariance_excess}.

\textbf{What the three equations do not enforce.}
For an isotropic velocity covariance, a realizable set of moments satisfies
\begin{equation}
 \Pi_f^2\leq \langle v_r'^2\rangle_f\Phi_f
             =\frac23 e_{t,f}\Phi_f .
 \label{eq:realizability}
\end{equation}
The local closure implemented here does not guarantee this bound under
all driving histories. Nonnegative $\Phi$ alone is insufficient.
The coupled Pi remap and kinetic energy transport without moment
transport do not establish it either.
The profile diagnostic reports
\begin{equation}
 \code{Pi_covariance_excess}=\Pi_f^2-\frac23e_{t,f}\Phi_f ;
\end{equation}
positive values indicate violation. This diagnostic imposes no flux cap.

The local entropy closure follows Flaskamp and Braun et al.~\cite{flaskamp,braun};
Kupka et al.~\cite{kupka} describe the nonlocal framework.
RSP2 retains kinetic energy transport but omits entropy moment transport.
Moment memory permits delayed or countergradient heat flux. It does not
guarantee numerical stability, covariance realizability or a stellar calibration.

\clearpage
\section{Source map}
Paths below are relative to the MESA repository. This map identifies the
implementation of the equations and state transitions described above.
\begin{center}
\small
\renewcommand{\arraystretch}{1.15}
\begin{tabularx}{\linewidth}{@{}p{0.28\linewidth} X@{}}
\toprule
Responsibility & Source and principal entry points\\
\midrule
Controls and constants &
\path{star/defaults/controls_dev.defaults};
\path{star/private/hydro_rsp2.f90}: \code{x_ALFAPI}, \code{x_ALFAPHI}.\\
Energy and moment rows &
\path{star/private/hydro_rsp2.f90}:
\code{do1_turbulent_energy_eqn}, \code{do1_rsp2_moment_eqns},
\code{rsp2_moment_rhs}, \code{rsp2_moment_source}.\\
Face EOS and gradients &
\path{star/private/hydro_gradient_support.f90}:
\code{get_rsp2_face_eos}, \code{get_rsp2_thermal_gradient},
\code{get_rsp2_Lrad_coeff}, \code{get_dPrad_dm_factors}.\\
Hydro temperature and energy &
\path{star/private/hydro_temperature.f90};
\path{star/private/hydro_energy.f90}.\\
Lengths and viscosity &
\path{star/private/tdc_hydro.f90}: public face and cell length functions;
\path{star/private/hydro_rsp2.f90}: \code{compute_Chi_cell},
\code{compute_Chi_face}, \code{compute_Eq_cell},
\code{compute_Uq_face}, \code{compute_Uq_dm_cell}.\\
Scales and trial state &
\path{star/private/hydro_rsp2.f90}: \code{set_etrb_start_vars},
\code{RSP2_adjust_vars_before_call_solver};
\path{star/private/solver_support.f90}.\\
Mode conversion &
\path{star/private/set_flags.f90}: \code{set_RSP2_3equation_flag};
\path{star/private/hydro_rsp2.f90}: \code{init_rsp2_moments}.\\
Mesh and saved state &
\code{star/private/hydro_rsp2.f90}: \code{remesh_rsp2_moments};
\path{star/private/mesh_adjust.f90},
\path{star/private/adjust_mesh_split_merge.f90},
\path{star/private/struct_burn_mix.f90},
\path{star/private/read_model.f90}, \path{star/private/write_model.f90},
\path{star/private/photo_in.f90}, \path{star/private/photo_out.f90}.\\
LNA and output &
\path{star/private/star_LNA_support.f90},
\path{star/private/star_LNA_turbulence_closures.f90},
\path{star/private/profile_getval.f90}.\\
\bottomrule
\end{tabularx}
\end{center}
Equation dispatch is in \path{star/private/hydro_eqns.f90};
the stored fields and layout are in \path{star_data/public/}.
No independent RSP3 state store bypasses the normal MESA solver.

\begingroup
\small
\begin{thebibliography}{9}
\bibitem{braun}
Braun, T. A. M., Ahlborn, F., Kupka, F. and Weiss, A. (2026).
\emph{Testing the 3-equation Kuhfuss Convection Model using the Sun}.
\href{https://arxiv.org/html/2604.06151v1}{arXiv:2604.06151v1}.
Section 2 gives local and nonlocal closures and their normalization;
Table 1 gives the solar calibrations.

\bibitem{flaskamp}
Flaskamp, M. (2003). \emph{Nichtlokale und zeitabh\"angige
Konvektion in Sternen}.
Doctoral thesis, Technische Universit\"at M\"unchen.
\href{https://mediatum.ub.tum.de/doc/602965/602965.pdf}{Thesis}.
Chapter 3 develops the entropy moment formulation.

\bibitem{kupka}
Kupka, F., Ahlborn, F. and Weiss, A. (2022).
\href{https://doi.org/10.1051/0004-6361/202243125}
{Astronomy \& Astrophysics, doi:10.1051/0004-6361/202243125}.
Appendix A gives the three equation Kuhfuss formulation.
\end{thebibliography}
\endgroup
\end{document}
